\cvsection{Expériences professionnelles}

\begin{cventries}
  \cventry
    {Stagiaire ingénieur de recherche en IA}
    {INRIA}
    {Bordeaux, France}
    {Mars 2026 -- Septembre 2026}
    {
      \begin{cvitems}
        \item {Analyse du besoin des chercheurs et développement en autonomie d’une application fonctionnelle intégrant un agent IA multi-étapes pour la recherche d’articles scientifiques.}
        \item {Développement de l’agent sans framework d’orchestration, avec intégration de l’API OpenAI, conception des prompts et gestion des appels d’outils, de l’état et de la mémoire.}
        \item {Mise en œuvre d’un pipeline RAG avec chunking, génération d’embeddings, indexation dans ChromaDB et recherche sémantique.}
        \item {Développement d’API REST en Python/FastAPI documentées avec OpenAPI, stockage dans MongoDB et interface React/Next.js permettant de choisir les messages transmis au modèle.}
        \item {Déploiement de LLM à poids ouverts avec vLLM sur un cluster HPC équipé de GPU NVIDIA et gestion de modèles locaux avec LM Studio sur Mac Studio.}
        \item {Comparaison de LLM et de quantifications avec Slurm, selon la réussite de la recherche, la latence d’inférence et la taille du contexte.}
        \item {Évaluation d’OpenClaw dans un environnement Docker isolé pour étudier les risques liés à l’accès des agents aux outils et aux fichiers.}
        \item {Conteneurisation de composants avec Docker et mise en place d’un pipeline CI/CD sur GitHub.}
      \end{cvitems}
    }

  \cventry
    {Stagiaire ingénieur de recherche en apprentissage par renforcement}
    {Université Karunya}
    {Coimbatore, Inde}
    {Juin 2025 -- Sept. 2025}
    {
      \begin{cvitems}
        \item {Développement d’un modèle d’apprentissage par renforcement et de son environnement d’entraînement pour la navigation autonome d’un robot hospitalier, puis déploiement et tests sur Raspberry Pi.}
      \end{cvitems}
    }
\end{cventries}
